\section*{Bab \ref{ch:g6:measure} --- Keliling, Luas, Volume}

\begin{solution}{exo:g6:measure:1}
$3.5$ m $= 350$ cm; \quad $420$ mm $= 42$ cm; \quad
$0.8$ km $= 800$ m; \quad $25\,000$ m $= 25$ km.
\end{solution}

\begin{solution}{exo:g6:measure:2}
Persegi panjang: $2 \times (8 + 3.5) = 2 \times 11.5 = 23$ cm.
Persegi: $4 \times 6.2 = 24.8$ cm.
Segitiga: $5 + 7 + 9 = 21$ cm.
\end{solution}

\begin{solution}{exo:g6:measure:3}
Satu putaran: $2\pi \times 50 = 100\pi \approx 314$ m. Untuk $2$ km
$= 2000$ m: $2000 \div 314 \approx 6.4$, jadi diperlukan $7$ putaran
penuh ($6$ putaran hanya menempuh kira-kira $1\,885$ m).
\end{solution}

\begin{solution}{exo:g6:measure:4}
Persegi panjang: $9 \times 4 = 36$ cm$^2$. Persegi: $7^2 = 49$ m$^2$.
\omterm{def:g6:shapes:triangles}{Segitiga siku-siku}: $\frac{6 \times 10}{2} = 30$ cm$^2$.
Cakram: $\pi \times 3^2 = 9\pi \approx 28$ cm$^2$.
\end{solution}

\begin{solution}{exo:g6:measure:5}
$3$ m$^2 = 30\,000$ cm$^2$; \quad
$45\,000$ cm$^2 = 4.5$ m$^2$; \quad
$2.5$ L $= 2\,500$ cm$^3$; \quad
$4\,500$ L $= 4.5$ m$^3$.
\end{solution}

\begin{solution}{exo:g6:measure:6}
Pagarnya (\omterm{def:g6:measure:perimeter}{kelilingnya}): $2 \times (120 + 85) = 2 \times 205 = 410$ m.
Luasnya: $120 \times 85 = 10\,200$ m$^2$.
\end{solution}

\begin{solution}{exo:g6:measure:7}
\omterm{def:g5:solids:def}{Balok}: $5 \times 4 \times 10 = 200$ cm$^3$. \omterm{def:g5:solids:def}{Kubus}:
$3 \times 3 \times 3 = 27$ cm$^3$.
\end{solution}

\begin{solution}{exo:g6:measure:8}
\omterm{def:g6:measure:perimeter}{Keliling} $16$ cm berarti panjang $+$ lebar $= 8$ cm. Misalnya
$6 \times 2$ (luasnya $12$ cm$^2$) dan $5 \times 3$ (luasnya $15$
cm$^2$) --- atau persegi $4 \times 4$ (luasnya $16$ cm$^2$). Makin
dekat ke bentuk persegi, makin besar luasnya.
\end{solution}

\begin{solution}{exo:g6:measure:9}
Luasnya: $10 \times 2 + 2 \times 6 = 20 + 12 = 32$ cm$^2$.

\omterm{def:g6:measure:perimeter}{Kelilingnya} (batang tegaknya di tengah bawah batang mendatar): telusurilah,
\[
10 + 2 + 4 + 6 + 2 + 6 + 4 + 2 = 36 \text{ cm}
\]
(atas; ujung kanan batang mendatar; bawah kanan; kanan batang tegak;
alas batang tegak; kiri batang tegak; bawah kiri; ujung kiri batang mendatar).
\end{solution}

\begin{solution}{exo:g6:measure:10}
\emph{1.} $V = 25 \times 10 \times 2 = 500$ m$^3$.

\emph{2.} $500$ m$^3 = 500 \times 1000 = 500\,000$ L.

\emph{3.} $500\,000 \div 50\,000 = 10$ jam.
\end{solution}

\begin{solution}{exo:g6:measure:11}
Kebun: $20^2 = 400$ m$^2$. Kolam: $\pi \times 3^2 = 9\pi$ m$^2$.
Rumput: $400 - 9\pi \approx 400 - 28.3 \approx 372$ m$^2$.
\end{solution}

\begin{solution}{exo:g6:measure:12}
\emph{1.} Batang kecil: $15 \times 6 \times 1 = 90$ cm$^3$. Batang
besar: $30 \times 12 \times 2 = 720$ cm$^3$. \omterm{def:g6:measure:volume}{Volumenya} dikalikan
$720 \div 90 = 8$ --- melipatduakan ketiga ukurannya mengalikan
\omterm{def:g6:measure:volume}{volumenya} dengan $2 \times 2 \times 2 = 8$.

\emph{2.} Delapan kali cokelatnya untuk empat kali harganya: ya,
batang yang besar dua kali lebih menguntungkan per gram.
\end{solution}

\begin{solution}{pb:g6:measure:1}
\textbf{1.} $1 \times 9$: \omterm{def:g6:measure:perimeter}{kelilingnya} $2 \times (1 + 9) = 20$ m,
luasnya $9$ m$^2$. \ $2 \times 8$: \omterm{def:g6:measure:perimeter}{kelilingnya} $2 \times 10 = 20$ m,
luasnya $16$ m$^2$.

\textbf{2.} \omterm{def:g6:measure:perimeter}{Kelilingnya} dua kali (panjang $+$ lebar)
(\cref{ex:g6:measure:perimeter}), jadi panjang $+$ lebar
$= 20 \div 2 = 10$ m. Tabelnya:
\begin{center}
\renewcommand{\arraystretch}{1.2}
\begin{tabular}{c|ccccc}
kandang & $1 \times 9$ & $2 \times 8$ & $3 \times 7$ & $4 \times 6$
& $5 \times 5$ \\
\hline
\omterm{def:g6:measure:area}{luas} (m$^2$) & $9$ & $16$ & $21$ & $24$ & $25$ \\
\end{tabular}
\end{center}
Persegi $5 \times 5$ yang terbaik.

\textbf{3.} $4.5 \times 5.5 = 24.75$ dan $4.9 \times 5.1 =
24.99$: makin lama makin dekat ke $25$, tetapi tetap di bawahnya.
Dugaan: persegi itu mengalahkan \emph{setiap} persegi panjang yang
berkeliling $20$, termasuk yang sisinya desimal.

\textbf{4.} Dari tabelnya: \omterm{def:g2:mult:def}{hasil kali} dua bilangan yang berjumlah
$10$ paling besar ketika kedua bilangannya \emph{sama} ($5$ dan
$5$). Aturan umumnya: untuk \omterm{def:g1:addition:def}{jumlah} yang tetap, \omterm{def:g2:mult:def}{hasil kalinya} paling
besar untuk bagian yang sama --- makin tidak sama pasangannya, makin
kecil \omterm{def:g2:mult:def}{hasil kalinya}.

\textbf{5.} Dari persegi $5 \times 5$, buanglah baris paling atas ---
jalur berisi $5$ persegi satuan --- sehingga tersisa kandang
$5 \times 4$; lalu rekatkan satu kolom berisi $4$ persegi satuan pada
salah satu ujungnya, sehingga kandangnya menjadi $6 \times 4$. Lima
persegi diambil dan hanya empat yang kembali: jalur yang ditambahkan
terletak sepanjang sisi $4$, yang lebih pendek daripada sisi $5$ yang
ditutupi jalur yang dibuang tadi. Kerugian bersihnya: satu meter
persegi ($25 \to 24$). Langkah berikutnya, dari $6 \times 4$ ke
$7 \times 3$, menukar baris berisi $6$ dengan kolom berisi $3$ dan
kehilangan tiga lagi ($24 \to 21$): makin jauh kandangnya dari bentuk
persegi, makin panjang jalur yang ia lepaskan dan makin pendek jalur
yang ia terima kembali.

\textbf{6.} Dengan $w + L + w = 20$, yaitu $L = 20 - 2w$:
\begin{center}
\renewcommand{\arraystretch}{1.2}
\begin{tabular}{c|ccccccc}
$w$ & $1$ & $2$ & $3$ & $4$ & $5$ & $6$ & $7$ \\
\hline
$L$ & $18$ & $16$ & $14$ & $12$ & $10$ & $8$ & $6$ \\
\omterm{def:g6:measure:area}{luas} & $18$ & $32$ & $42$ & $48$ & $50$ & $48$ & $42$ \\
\end{tabular}
\end{center}
Pemenangnya adalah $5 \times 10$, dengan \omterm{def:g6:measure:area}{luas} $50$ m$^2$ --- dua kali
lebih panjang (sepanjang dinding) daripada lebarnya: \emph{bukan} persegi.

\textbf{7.} Cerminkan kandangnya terhadap dindingnya: kandangnya
beserta bayangan cerminnya membentuk kandang ganda yang memakai
pagarnya dua kali, yaitu $40$ m pagar mengelilinginya (sisi dindingnya
kini berada di dalam). Menurut Bagian I, persegi panjang terbaik yang
berkeliling $40$ adalah persegi $10 \times 10$, dengan \omterm{def:g6:measure:area}{luas} $100$
m$^2$. Kandang sebenarnya yang terbaik adalah setengah kandang ganda
terbaik itu: $10$ m sepanjang dinding, $5$ m ke dalam, luasnya $50$
m$^2$ --- tepat pemenang pada tabelnya, yang kini terjelaskan.

\textbf{8.} Lingkarnya $20 = 2 \times \pi \times r$, jadi
$r = 20 \div 6.28 \approx 3.18$ m. Luasnya:
$\pi r^2 \approx 3.14 \times 3.18 \times 3.18 \approx
31.8$ m$^2$ --- jauh lebih besar daripada $25$ m$^2$ milik persegi.
Dido tahu apa yang ia lakukan: untuk panjang tepi yang diberikan,
lingkaranlah yang melingkupi paling banyak.

\textbf{9.} \omterm{def:g6:measure:perimeter}{Kelilingnya}: $1 \times 36$: $74$ m; \ $2 \times 18$:
$40$ m; \ $3 \times 12$: $30$ m; \ $4 \times 9$: $26$ m; \
$6 \times 6$: $24$ m. Persegi lagi --- pagar paling sedikit untuk
\omterm{def:g6:measure:area}{luas} yang diberikan. Inilah Bagian I yang dibalik: \omterm{def:g6:measure:perimeter}{keliling} tetap
dengan \omterm{def:g6:measure:area}{luas} terbesar, atau \omterm{def:g6:measure:area}{luas} tetap dengan \omterm{def:g6:measure:perimeter}{keliling} terkecil;
keduanya menobatkan persegi (dan, di luar semua persegi panjang, lingkaran).

\textbf{10.} Dinding yang bulat adalah tepi terpendek untuk ruang yang
dilingkupinya (pertanyaan 8), jadi kaleng, pipa atau tangki yang bulat
memuat isi paling banyak dengan logam paling sedikit --- bahan itu uang.

\textbf{11.} Misalnya $50$ m $\times$ $0.01$ m: \omterm{def:g6:measure:perimeter}{kelilingnya}
$2 \times (50 + 0.01) = 100.02$ m, luasnya
$50 \times 0.01 = 0.5$ m$^2$. Panjang dan tipis: bertepi
berkilo-kilometer, tetapi tanahnya hampir tidak ada.

\textbf{12.} Salah. Kandang pada pertanyaan 11 berkeliling
$100.02$ m dan berluas $0.5$ m$^2$, sedangkan kandang $5 \times 5$
berkeliling $20$ m dan berluas $25$ m$^2$: \omterm{def:g6:measure:perimeter}{kelilingnya} lebih besar,
luasnya jauh lebih kecil.

\textbf{13.} Luasnya: satu persegi satuan hilang,
$24 - 1 = 23$ cm$^2$. \omterm{def:g6:measure:perimeter}{Kelilingnya}: ketika ditelusuri, takiknya
menggantikan $1$ cm dinding lurus dengan tiga sisi persegi kecil,
$1 + 1 + 1 = 3$ cm: \omterm{def:g6:measure:perimeter}{kelilingnya} bertambah dari $20$ cm menjadi
$20 - 1 + 3 = 22$ cm. Membuang bahan justru \emph{memanjangkan}
tepinya.

\textbf{14.} Setiap perbesaran pada pantai yang sesungguhnya
memunculkan teluk, batu dan cerukan yang baru --- takik di atas takik,
dan pertanyaan 13 memperlihatkan bahwa setiap takik menambah panjang
tepi sambil hampir tidak mengubah luasnya. Jadi panjang yang terukur
terus bertambah seiring tingkat kerinciannya (``paradoks garis
pantai''), sedangkan \omterm{def:g6:measure:area}{luas} yang terukur mengendap: \omterm{def:g6:measure:perimeter}{keliling} dan \omterm{def:g6:measure:area}{luas}
benar-benar besaran yang saling bebas.

\textbf{15.} Alasan cerminnya: kandang gandanya akan memakai
$2 \times 36 = 72$ m pagar, dan persegi panjang terbaik yang
berkeliling $72$ adalah persegi $18 \times 18$. Jadi kandang terbaiknya
$18$ m sepanjang dinding dan $9$ m ke dalam: periksa pagarnya
$9 + 18 + 9 = 36$ m, luasnya $18 \times 9 = 162$ m$^2$. Di lapangan
terbuka, yang terbaik adalah persegi $9 \times 9$, luasnya $81$ m$^2$.
Dinding lumbung itu tepat \emph{melipatduakan} tanah sang petani.

\end{solution}
